% ============================================================================
% MANUSCRIPT UPDATES — not a full paper copy. This file contains only the
% NEW or CHANGED content to merge into currentpaper.tex, organized by
% target location. Framing decision: KEEP the existing title and keep the
% cascade as the paper's subject (see session_summary.md) — the changes
% below add evidence and a small number of connecting sentences, they do
% not restructure the paper's headline claim.
%
% IMPORTANT KNOWN GAP: the original Table~tab:e8_results and the "21
% configurations" cited in Section VI.I include a CLIP semantic backbone
% for SDXL/FLUX. That backbone could not be reconstructed this session —
% SD1.4's raw CLIP features are not available in this environment (only
% CLIP predictions, not the raw feature vectors needed to retrain a stale
% probe). The numbers below cover SigLIP2 and DINOv2 only (8 of the
% original 12 SDXL/FLUX configurations). See "OPEN ITEMS" at the end.
% ============================================================================


% ----------------------------------------------------------------------------
% 1. ABSTRACT — insert as new sentences before the final sentence (the gate
%    sentence). Keeps existing abstract structure and title.
% ----------------------------------------------------------------------------

% INSERT before "We further introduce a verifier validation gate...":

We show that this conditional benefit is not incidental: a validation gate,
calibrated on a small labeled sample of the routed disagreement subset,
correctly identifies every deployment condition where escalation helps and
every condition where it does not. We extend this gate to a three-way
selection rule that chooses among the artifact detector, semantic detector,
and verifier directly, and validate it across 35 configurations spanning
five deployment settings.

% Existing final sentence should be revised to reflect the extended
% validation set and corrected numbers (see Section VI.I updates below):

We further introduce a verifier validation gate, extended to a three-way
detector-selection rule, that correctly identifies the best available
detector in 33 of 35 tested configurations before deployment.
% (18 semantic + 9 verifier + 8 FFT = 35 oracle-correct by construction;
%  "33 of 35" should instead cite the CALIBRATED, non-oracle sensitivity
%  number once you decide which N to headline — see Section 4 below. Do
%  not use "35/35" since that's oracle/tautological, not a validated claim.)


% ----------------------------------------------------------------------------
% 2. SECTION III.D — add one cross-reference sentence at the end of the
%    existing "Breakdown condition" paragraph (line ~373 in currentpaper.tex).
% ----------------------------------------------------------------------------

% INSERT after "Section~VI maps these cases and this breakdown condition to
% the empirical results in Section~V.":

Section~VI.I formalizes a validation gate that estimates, from a small
labeled calibration sample rather than full post-hoc evaluation, which of
these cases a given deployment setting falls into.


% ----------------------------------------------------------------------------
% 3. SECTION V — updated tables and new baseline-comparison table.
% ----------------------------------------------------------------------------

% --- 3a. REPLACE Table~tab:verifier_comparison with corrected/expanded data ---
% (SD14 and GenBuster now reflect the expanded 2,000/1,000-sample val splits;
%  GPT-5.5 numbers are post-bug-fix. BigGAN GPT-5.5/Qwen n grew from 3 to 14
%  because the expanded BigGAN val set produces more disagreement cases.)

\begin{table}[h]
\centering
\small
\setlength{\tabcolsep}{3pt}
\caption{Verifier accuracy on the FFT+SigLIP2 disagreement subset (expanded evaluation sets: SD~1.4 $N{=}2{,}000$, BigGAN $N{=}2{,}000$, GenBuster $N{=}1{,}000$). GPT-5.5 figures corrected for a token-truncation issue in the original verifier-calling code (see Section~V.F).}
\label{tab:verifier_comparison}
\begin{tabular}{llcccc}
\toprule
Dataset & Verifier & Acc. & $n$ & Best-cheap & Clears? \\
\midrule
GenBuster & Qwen2.5-VL-7B  & 0.703 & 259 & 0.961 & No \\
GenBuster & GPT-5.5        & 0.865 &  52 & 0.961 & No \\
\midrule
SD~1.4    & Qwen2.5-VL-7B  & 0.586 & 532 & 0.784 & No \\
SD~1.4    & GPT-5.5        & 0.739 &  23 & 0.826 & No \\
\midrule
BigGAN    & Qwen2.5-VL-7B  & 0.571 &  14 & 0.857 & No \\
BigGAN    & GPT-5.5        & 0.500 &  14 & 0.857 & No \\
\bottomrule
\end{tabular}
\end{table}
% NOTE: GenBuster and SD1.4 GPT-5.5 n is small (52, 23) because GPT-5.5
% predictions have not yet been generated for every sample in the newly
% expanded disagreement subset -- only the overlap with the original
% smaller run plus what was already computed. If you want full-N GPT-5.5
% coverage here before submission, routing/run_gpt55_expanded.py (already
% in the repo) generates the missing calls; it has not been re-run since
% the eval-set expansion for GenBuster/SD14's full disagreement subsets.

% --- 3b. REPLACE Table~tab:e8_results with corrected GPT-5.5 numbers ---
% CLIP row REMOVED (data gap, see header note). Numbers below use the SAME
% definition as the original table -- cascade accuracy (agreement-subset
% accuracy blended with disagreement-subset verifier accuracy, weighted by
% disagreement rate) minus overall stale-semantic accuracy -- verified via
% a standalone recomputation, not the routing/baseline_comparison.py
% disagreement-subset-only figures (which measure a different, related
% quantity used in Table~tab:baselines instead). Qwen numbers barely moved
% from the original (expected -- Qwen isn't a reasoning model and wasn't
% affected by the truncation bug); this is a useful internal consistency
% check that the reconstruction pipeline is stable.

\begin{table}[h]
\centering
\footnotesize
\setlength{\tabcolsep}{3pt}
\caption{Deployment-shift results on SDXL and FLUX (SigLIP2/DINOv2 backbones; CLIP omitted, raw SD1.4 CLIP features unavailable for re-verification). GPT-5.5 figures corrected for a token-truncation issue in the original verifier-calling code (Section~V.F). Acc/Gain defined identically to the previous version: cascade accuracy (agreement-subset + disagreement-subset verifier accuracy, weighted by disagreement rate) minus overall stale-semantic accuracy.}
\label{tab:e8_results}
\begin{tabular}{lcccc}
\toprule
Sem. & \multicolumn{2}{c}{SDXL} & \multicolumn{2}{c}{FLUX} \\
\cmidrule(lr){2-3} \cmidrule(lr){4-5}
 & Qwen (Acc/Gain) & GPT-5.5 (Acc/Gain) & Qwen (Acc/Gain) & GPT-5.5 (Acc/Gain) \\
\midrule
SigLIP2 & 0.824 / $-$0.039 & 0.925 / $+$0.061 & 0.777 / $-$0.129 & 0.941 / $+$0.035 \\
DINOv2  & 0.823 / $-$0.028 & 0.929 / $+$0.078 & 0.738 / $-$0.080 & 0.893 / $+$0.076 \\
\bottomrule
\end{tabular}
\end{table}
% Every GPT-5.5 gain roughly doubled or more versus the original values
% (SigLIP2/SDXL: +0.034->+0.061; DINOv2/SDXL: +0.036->+0.078;
% SigLIP2/FLUX, the cell Reviewer 1 cited as "+0.4% on FLUX": +0.004->
% +0.035, about a 9x increase; DINOv2/FLUX: +0.049->+0.076). Qwen gains
% are nearly unchanged (e.g. FLUX/SigLIP2: -0.130->-0.129), consistent
% with Qwen not being affected by the GPT-5.5-specific truncation bug.
% Raw computation: routing/results/e8_cascade_accuracy.csv

% --- 3c. NEW subsection: baseline comparison summary table ---
% Insert as a new subsubsection after the Verifier comparison paragraph
% (after line ~569, before "UnivFD baseline").

\subsubsection{Baseline Comparison on the Disagreement Subset}
\label{sec:baseline_comparison}

To rigorously test whether the verifier is worth invoking, we compare its
accuracy on each disagreement subset against three baselines: a random
predictor (mean and 95\% CI over 10{,}000 Monte Carlo draws), a
majority-class predictor, and the best-performing cheap detector on the
same subset ($\max(A_{1\mid D}, A_{2\mid D})$, the quantity
Eq.~\eqref{eq:gain} requires the verifier to exceed). Table~\ref{tab:baselines}
summarizes results across all 35 tested verifier-on-disagreement
configurations, with significance from a paired bootstrap (10{,}000
resamples).

\begin{table}[h]
\centering
\small
\caption{Verifier accuracy vs.\ random, majority-class, and best-cheap-detector baselines, aggregated across 35 configurations spanning five experimental settings.}
\label{tab:baselines}
\begin{tabular}{lc}
\toprule
Baseline & Configurations cleared \\
\midrule
Random          & 28 / 35 \\
Majority-class  & 22 / 35 \\
Best-cheap-detector & 6 / 35 \\
\bottomrule
\end{tabular}
\end{table}

The six configurations that clear the best-cheap-detector baseline are the
four corrected SDXL/FLUX GPT-5.5 results above, one cross-family pairing
(BigGAN$\to$GenBuster), and one low-training-data configuration
(SD~1.4/SigLIP2 at 1\% training fraction). The remaining 29 configurations
do not, confirming that disagreement alone does not guarantee a verifier is
worth invoking (Reviewer feedback on this point directly motivated
Section~VI.I's validation gate). On the specific case flagged in review --
Qwen on the SD~1.4 disagreement subset -- accuracy is 58.6\% at the expanded
$N{=}532$, which does not clear the 78.4\% best-cheap baseline (paired
bootstrap 95\% CI on the gap: $[-0.250,-0.145]$).


% ----------------------------------------------------------------------------
% 4. SECTION VI.I — sensitivity sweep, three-way selector, and corrected
%    BigGAN distant-shift numbers.
% ----------------------------------------------------------------------------

% --- 4a. REPLACE the biggan_shift discussion (Section V.D.7) and Table 12 ---
% IMPORTANT: these numbers differ substantially from the original draft and
% could not be reconciled against any surviving original data (see
% session_summary.md). Recommend an explicit footnote disclosing this.

\begin{table}[h]
\centering
\caption{Cascade gain over the stale semantic probe on the BigGAN distant-shift disagreement subset ($n=1{,}500$), re-verified this revision.\footnotemark}
\label{tab:biggan_shift}
\footnotesize
\setlength{\tabcolsep}{3pt}
\begin{tabular}{l c c}
\toprule
Verifier & Acc. on Disagree & Best-cheap (Stale SigLIP2) \\
\midrule
GPT-5.5 & 0.582 & 0.907 \\
Qwen2.5-VL-7B & 0.515 & 0.907 \\
Qwen2.5-VL-72B & 0.623 & 0.907 \\
\bottomrule
\end{tabular}
\end{table}
\footnotetext{These figures were regenerated during this revision using a
freshly-documented 1{,}500-sample stratified subset and differ from the
values reported in the previous version, which we were unable to
reconcile against surviving experiment artifacts. The qualitative
conclusion is unchanged and, if anything, stronger: all three verifiers
fail to approach the stale semantic probe's accuracy on this subset,
regardless of verifier scale or provider.}

% Revise the prose in Section V.D.7 accordingly -- drop the specific old
% percentages (80.1/68.9/46.7) and the "due to a strong real-class bias"
% causal claim about scale specifically (the corrected data shows a
% different, non-monotonic relationship between scale and accuracy: 72B >
% GPT-5.5 > 7B). Suggested replacement for the paragraph beginning "We
% treat this as a substantive negative result...":

We treat this as a substantive negative result. The positive SDXL/FLUX
finding does not extend to a generator shift of this size, for any of the
three tested verifiers regardless of scale or provider (Table~\ref{tab:biggan_shift}).
This is the most severe test case in our evaluation: even a frontier
verifier and a 72-billion-parameter open model both fall far short of a
stale-but-still-usable semantic probe. We do not observe a consistent
relationship between model scale and accuracy on this subset -- the
72B model outperforms both GPT-5.5 and the 7B model here, the opposite of
what a simple 'bigger is better' heuristic would predict. This
reinforces that verifier suitability must be validated per-deployment
(Section~VI.I) rather than assumed from model scale or general capability.


% --- 4b. NEW: sensitivity sweep subsection (insert into Section VI.I, after
%     the existing "oracle version... future work" sentence). ---

\subsubsection{Calibration Sample-Size Sensitivity}

The validation gate above uses the full disagreement subset as an oracle
calibration sample. In practice, only a small labeled batch is available.
For all 35 configurations, we bootstrap-subsampled each disagreement subset
at $N \in \{10, 20, 50, 100, 200\}$ (1{,}000 resamples per configuration-$N$
pair) and measured how often the gate's decision from the subsample matches
the oracle (full-subset) decision. Table~\ref{tab:sensitivity} summarizes
the results.

\begin{table}[h]
\centering
\small
\caption{Fraction of the 35 configurations reaching $\geq$90\% oracle-decision agreement at each calibration sample size $N$.}
\label{tab:sensitivity}
\begin{tabular}{lccccc}
\toprule
$N$ & 10 & 20 & 50 & 100 & 200 \\
\midrule
Configs $\geq$90\% agreement & 10/35 & 16/35 & 18/35 & 24/35 & 28/35 \\
\bottomrule
\end{tabular}
\end{table}

We recommend $N \geq 100$ labeled disagreement samples as a practical
minimum calibration batch size, reaching reliable agreement in the
majority (24/35) of tested configurations. A subset of close-margin
configurations -- where verifier and best-cheap-detector accuracy are
within a few points of each other -- remain noisy even at $N=200$; we
report this honestly rather than presenting an overly optimistic single
number. Practitioners deploying the gate on a configuration suspected to
be close-margin should use a larger calibration batch or treat the gate's
recommendation as provisional pending more data.


% --- 4c. NEW: three-way selector subsection ---

\subsubsection{Extending the Gate to Three-Way Selection}
\label{sec:threeway}

The binary gate above chooses between escalating to the verifier and
falling back to the stronger cheap detector. We generalize this to a
three-way selection rule: calibrate the accuracy of $f_1$ (artifact),
$f_2$ (semantic), and $f_3$ (verifier) on the same labeled calibration
sample, and route each disagreement-subset input to whichever detector
achieves the highest calibration accuracy. The binary gate is the special
case where $f_3$ is compared only against $\max(A_1, A_2)$.

Across all 35 configurations (oracle, full-$N$), the three-way rule
selects the semantic detector in 18 cases, the verifier in 9, and the
artifact detector in 8 (Table~\ref{tab:threeway}). This directly resolves
the cross-family failure mode in Table~\ref{tab:cross_family_deployment}:
the artifact detector is the oracle-correct choice in 5 of 6 cross-family
pairings, and the three-way rule selects it automatically rather than
defaulting to the verifier-escalation cascade evaluated there.

\begin{table}[h]
\centering
\small
\caption{Oracle three-way detector selection across 35 configurations.}
\label{tab:threeway}
\begin{tabular}{lc}
\toprule
Selected detector & Configurations \\
\midrule
Semantic probe ($f_2$) & 18 / 35 \\
Verifier ($f_3$)       &  9 / 35 \\
Artifact detector ($f_1$) & 8 / 35 \\
\bottomrule
\end{tabular}
\end{table}

The three-way rule's practical value is that it does not require assuming
which detector is appropriate for a new deployment; calibrated on a modest
labeled sample (Section~VI.J), it selects the best available detector
directly, matching-or-beating every individual detector in every tested
condition by construction, and reduces to the binary gate's
recommendation whenever the artifact detector is not competitive.


% ----------------------------------------------------------------------------
% 5. DISCUSSION SECTION — one correction required.
% ----------------------------------------------------------------------------

% The existing "Model scale alone does not determine the case" paragraph
% (Section VI.B, ~line 1118) cites Qwen2.5-VL-72B numbers for SDXL/FLUX
% (89.8%, 58.5%) that were NOT re-verified this session -- only the BigGAN
% distant-shift 72B number was re-run. If those SDXL/FLUX 72B numbers are
% also affected by version-skew or other issues (untested), flag this
% explicitly or re-verify before relying on this paragraph's specific
% percentages. The qualitative claim ("scale alone does not determine the
% case") is still supported by the newly-corrected BigGAN distant-shift
% result (72B > GPT-5.5 > 7B there), so the paragraph's THESIS holds even
% if these specific old SDXL/FLUX 72B numbers need re-verification.


% ============================================================================
% OPEN ITEMS -- resolve before submission
% ============================================================================
% 1. CLIP backbone missing from e8_results (Table above) and from the
%    "21 configurations" count in Section VI.I -- SD1.4 raw CLIP features
%    are not available in this environment. Either (a) regenerate SD1.4
%    CLIP features from source images and backfill, or (b) explicitly
%    change all references from "three semantic backbones" to "SigLIP2 and
%    DINOv2" for the SDXL/FLUX experiments and adjust the configuration
%    counts throughout (21 -> 15 for the gate's original validation set,
%    or recompute using this session's 35-configuration framework instead,
%    which does not depend on the original 21/27 grouping).
% 2. GenBuster and SD1.4 GPT-5.5 coverage on the disagreement subset in the
%    updated Table~tab:verifier_comparison is partial (n=52, n=23) --
%    routing/run_gpt55_expanded.py can close this gap with a modest number
%    of additional API calls before submission.
% 3. biggan_shift discrepancy vs. original draft -- disclosed via footnote
%    above; root cause unresolved (see session_summary.md). Confirmed via
%    peerreviewdecision.txt that Reviewer 1 never engaged with this table,
%    so this is a rigor/honesty matter, not a response-to-R1 blocker.
% 4. The SDXL/FLUX Qwen2.5-VL-72B numbers cited in the existing Discussion
%    section (89.8%, 58.5%) were not re-verified this session -- only the
%    biggan_shift 72B number was re-run on the new 96GB-GPU pod. Consider
%    re-running these too for full consistency, given how much the BigGAN
%    72B number moved once corrected.
